\chapter{Analisi dei requisiti}
\label{cap:analisi-requisiti}

\intro{Il presente capitolo descrive le funzionalità del prodotto, formalizzate sotto forma di requisiti di progetto, e relativi casi d’uso associati.}\\

\section{Analisi del sistema}
Il prodotto sviluppato durante il periodo di stage mette a disposizione dello sviluppatore un sistema che permette la creazione automatizzata di test e il monitoraggio dello stato del sistema \emph{Sia.Core} tramite dashboard.
\\
Dal punto di vista funzionale, la soluzione si compone di due macro-sottosistemi:
\begin{itemize}
  \item un sottosistema di \emph{testing} automatico, integrato nel \emph{framework} \emph{Sia.Core} e articolato su \emph{unit test} e \emph{integration test} lato \emph{server} (logica di business, attendibilità dei dati nel database, procedure SQL e prospettive di configurazione) e su \emph{unit test} e test \emph{end-to-end} lato \emph{client} (logica applicativa e validazione dell'esperienza utente);
  \item un sottosistema di \emph{osservabilità}, basato sullo standard \gls{otel} e sullo stack \emph{Grafana LGTM} (\emph{Loki}, \emph{Grafana}, \emph{Tempo}, \emph{Mimir/Prometheus}), che raccoglie \emph{trace}, metriche e \emph{log} prodotti dall'applicazione e ne permette la visualizzazione e la correlazione tramite dashboard.
\end{itemize}
Trasversalmente ai due sottosistemi, è stato adottato un agente di intelligenza artificiale come assistente allo sviluppo in modo strutturato, con la creazione di agenti specializzati, di \emph{skill} procedurali e l'integrazione del \gls{mcp} per consentire all'agente di interrogare direttamente \emph{Grafana}, il database e gli strumenti di test.
I due sottosistemi sono indipendenti sul piano tecnologico ma complementari sul piano metodologico: la telemetria fornisce il \emph{feedback} che guida il miglioramento iterativo dei test, mentre i test supportano la validazione continua della strumentazione.

\section{Casi d'uso}
La presente sezione illustra i casi d’uso analizzati, corredati dai relativi diagrammi.
I diagrammi dei casi d’uso (in inglese \emph{Use Case Diagram}) costituiscono una tipologia di diagrammi \gls{uml} finalizzati alla rappresentazione delle funzionalità e dei servizi offerti da un sistema, dal punto di vista degli attori che interagiscono con esso.

\subsection{Attori}
Gli attori nei casi d’uso rappresentano entità esterne al sistema, quali utenti o altri servizi, che interagiscono con esso.
In merito al sistema sviluppato si possono individuare come attori:
\begin{itemize}
  \item \textbf{Sviluppatore}: è l'attore principale che utilizza il sistema, sia per generare e revisionare i test sia per consultare le dashboard di telemetria durante le attività di diagnosi e \emph{tuning}.
  \item \textbf{Agente AI}: agente di codifica basato su un modello linguistico di grandi dimensioni (\gls{llm}), invocato dall'ambiente di sviluppo come \emph{pair engineer}. Opera tramite \emph{skill} procedurali e \gls{mcp} per generare test, propagare la telemetria ai moduli del \emph{framework} e interrogare \emph{Grafana}, database e strumenti di test.
  \item \textbf{Sia.Core}: applicazione \emph{full-stack} oggetto dell'osservazione e del \emph{testing}, sorgente dei segnali di telemetria e obiettivo della suite di test.
\end{itemize}

\subsection{Elenco casi d'uso}

\begin{itemize}
  \item UC1 --- Generazione assistita di \emph{unit test} backend tramite agente AI.
        \begin{itemize}
          \item UC1.1 --- Identificazione del componente da coprire.
          \item UC1.2 --- Invocazione dell'agente AI tramite \emph{skill} dedicata.
          \item UC1.3 --- Analisi del codice e generazione dei test.
          \item UC1.4 --- Revisione e integrazione nella suite.
        \end{itemize}
  \item UC2 --- Generazione assistita di \emph{integration test} su database.
        \begin{itemize}
          \item UC2.1 --- Selezione dell'oggetto da validare.
          \item UC2.2 --- Interrogazione del database e delle prospettive via \gls{mcp}.
          \item UC2.3 --- Generazione dell'\emph{integration test}.
          \item UC2.4 --- Revisione e integrazione nella suite.
        \end{itemize}
  \item UC3 --- Generazione assistita di \emph{unit test} e test \emph{end-to-end} \emph{frontend}.
        \begin{itemize}
          \item UC3.1 --- Identificazione del componente o del flusso utente da coprire.
          \item UC3.2 --- Generazione di \emph{unit test} per la logica del componente.
          \item UC3.3 --- Generazione del test \emph{end-to-end}.
          \item UC3.4 --- Revisione e integrazione nella suite.
        \end{itemize}
  \item UC4 --- Esecuzione della suite di test e produzione del report di copertura.
        \begin{itemize}
          \item UC4.1 --- Avvio della suite di test.
          \item UC4.2 --- Esecuzione dei test integrata nel processo di \emph{build}.
          \item UC4.3 --- Generazione del report di copertura del codice.
          \item UC4.4 --- Gestione dei test falliti (estensione).
        \end{itemize}
  \item UC5 --- Propagazione della telemetria a un nuovo modulo del \emph{framework} tramite \emph{skill} dedicate.
        \begin{itemize}
          \item UC5.1 --- Identificazione del modulo da strumentare.
          \item UC5.2 --- Invocazione della \emph{skill} di propagazione della telemetria.
          \item UC5.3 --- Introduzione della classe di diagnostica e registrazione della sorgente nel \emph{pipeline} \gls{otel}.
          \item UC5.4 --- Definizione di metriche con cardinalità controllata.
          \item UC5.5 --- Validazione dei segnali sullo stack \emph{Grafana LGTM}.
        \end{itemize}
  \item UC6 --- Visualizzazione delle metriche, \emph{trace} e \emph{log} di un modulo tramite dashboard \emph{Grafana LGTM}.
        \begin{itemize}
          \item UC6.1 --- Apertura della dashboard del modulo.
          \item UC6.2 --- Selezione delle variabili \emph{template}.
          \item UC6.3 --- Consultazione correlata di metriche, \emph{trace} e \emph{log}.
        \end{itemize}
  \item UC7 --- Diagnosi di una richiesta lenta a partire da un pannello percentili, tramite \emph{exemplar} verso la \emph{trace} e navigazione ai \emph{log} correlati.
        \begin{itemize}
          \item UC7.1 --- Individuazione del campione anomalo sul pannello percentili.
          \item UC7.2 --- Navigazione all'\emph{exemplar} associato al campione.
          \item UC7.3 --- Apertura della \emph{trace} corrispondente in \emph{Tempo}.
          \item UC7.4 --- Navigazione ai \emph{log} correlati tramite \texttt{TraceId}.
        \end{itemize}
  \item UC8 --- Identificazione dell'utente coinvolto in una \emph{trace} tramite l'arricchimento automatico degli \emph{span}.
        \begin{itemize}
          \item UC8.1 --- Apertura della \emph{trace} di interesse.
          \item UC8.2 --- Lettura degli attributi \texttt{enduser.id} e \texttt{enduser.name} dello \emph{span} radice.
          \item UC8.3 --- Filtro di ulteriori \emph{trace} e \emph{log} per lo stesso utente.
        \end{itemize}
  \item UC9 --- Creazione o modifica di una dashboard \emph{Grafana} tramite agente AI e \gls{mcp}.
        \begin{itemize}
          \item UC9.1 --- Specifica delle esigenze di \emph{monitoring} all'agente AI.
          \item UC9.2 --- Validazione delle query PromQL/TraceQL/LogQL via \gls{mcp}.
          \item UC9.3 --- Pubblicazione o aggiornamento della dashboard sull'istanza \emph{Grafana}.
          \item UC9.4 --- Gestione di una query non valida (estensione).
        \end{itemize}
\end{itemize}
\newpage
\begin{usecase}{1}{Generazione assistita di unit test backend}
  \begin{figure}[!h]
    \centering
    \includegraphics[width=0.9\columnwidth]{usecase/UC1}
    \caption{UC1: Generazione assistita di unit test backend}
  \end{figure}

  \usecaseactors{Sviluppatore, Agente AI}
  \usecasepre{L'ambiente di sviluppo del \emph{framework} \emph{Sia.Core} è aperto, l'agente AI è configurato con la \emph{skill} dedicata agli \emph{unit test} \emph{backend} e nel progetto è presente il \emph{test runner} \emph{.NET}}
  \usecasedesc{Lo sviluppatore vuole introdurre \emph{unit test} su un servizio o un \emph{repository} del \emph{framework backend} privo di copertura adeguata, avvalendosi dell'agente AI}
  \usecasepost{Nuovi \emph{unit test} sono presenti nel progetto, eseguibili nella \emph{build}, e contribuiscono all'aumento della copertura del codice}
  \usecasescenario{
    \begin{enumerate}
      \item Lo sviluppatore identifica il componente da coprire (UC1.1);
      \item Lo sviluppatore invoca l'agente AI tramite \emph{prompt} (UC1.2);
      \item L'agente analizza il codice del componente e genera gli \emph{unit test} (UC1.3);
      \item Lo sviluppatore revisiona i test e li integra nella suite (UC1.4).
    \end{enumerate}
  }
  \label{uc:unit-test-backend}
\end{usecase}

\begin{usecase}{1.1}{Identificazione del componente da coprire}
  \usecasedesc{Lo sviluppatore individua un servizio o un \emph{repository} privo di copertura adeguata, anche basandosi sul report di copertura prodotto da una precedente esecuzione della suite}
  \label{uc:unit-test-backend-target}
\end{usecase}

\begin{usecase}{1.2}{Invocazione dell'agente AI tramite prompt}
  \usecasedesc{Lo sviluppatore richiama l'agente AI passando il riferimento al componente, l'agente riconosce il contesto del \emph{prompt} e recupera la \emph{skill} di generazione \emph{unit test}}
  \label{uc:unit-test-backend-invoke}
\end{usecase}

\begin{usecase}{1.3}{Analisi del codice e generazione dei test}
  \usecasedesc{L'agente analizza il codice del componente, ne deduce i comportamenti attesi e i casi limite, e produce \emph{unit test} che esercitano la logica di business}
  \label{uc:unit-test-backend-generate}
\end{usecase}

\begin{usecase}{1.4}{Revisione e integrazione nella suite}
  \usecasedesc{Lo sviluppatore valuta la correttezza e la pertinenza dei test, li adatta se necessario e li integra nella suite del progetto}
  \label{uc:unit-test-backend-review}
\end{usecase}

\begin{usecase}{2}{Generazione assistita di integration test su DB}
  \begin{figure}[!h]
    \centering
    \includegraphics[width=0.9\columnwidth]{usecase/UC2}
    \caption{UC2: Generazione assistita di integration test su DB}
  \end{figure}

  \usecaseactors{Sviluppatore, Agente AI}
  \usecasepre{Il \gls{mcp} è configurato verso il database del \emph{framework} e l'agente AI dispone delle \emph{skill} dedicate agli \emph{integration test}}
  \usecasedesc{Lo sviluppatore vuole verificare l'attendibilità dei dati nel database, il corretto comportamento di una procedura SQL o la coerenza di una prospettiva di configurazione di un componente, sfruttando l'agente AI e l'accesso diretto via \gls{mcp}}
  \usecasepost{La suite di \emph{integration test} verifica in modo automatizzato la consistenza dell'oggetto selezionato}
  \usecasescenario{
    \begin{enumerate}
      \item Lo sviluppatore seleziona l'oggetto da validare (UC2.1);
      \item L'agente AI ispeziona via \gls{mcp} schema, dati e procedure coinvolte (UC2.2);
      \item L'agente genera l'\emph{integration test} (UC2.3);
      \item Lo sviluppatore revisiona e integra i test nella suite (UC2.4).
    \end{enumerate}
  }
  \label{uc:integration-test-backend}
\end{usecase}

\begin{usecase}{2.1}{Selezione dell'oggetto da validare}
  \usecasedesc{Lo sviluppatore identifica l'oggetto target (tabella, vista, \emph{stored procedure}, prospettiva) e il criterio di validità che deve essere verificato}
  \label{uc:integration-test-target}
\end{usecase}

\begin{usecase}{2.2}{Interrogazione del database e delle prospettive via MCP}
  \usecasedesc{L'agente AI interroga il database e il filesystem delle prospettive per ricostruire schema, dati e dipendenze dell'oggetto target}
  \label{uc:integration-test-mcp}
\end{usecase}

\begin{usecase}{2.3}{Generazione dell'integration test}
  \usecasedesc{L'agente produce un \emph{integration test} che esegue le operazioni necessarie e ne verifica la conformità al criterio di validità formalizzato}
  \label{uc:integration-test-generate}
\end{usecase}

\begin{usecase}{2.4}{Revisione e integrazione nella suite}
  \usecasedesc{Lo sviluppatore valuta il test, lo adatta al contesto specifico e lo integra nella suite}
  \label{uc:integration-test-review}
\end{usecase}
\newpage
\begin{usecase}{3}{Generazione assistita di unit test e E2E frontend}
  \begin{figure}[!h]
    \centering
    \includegraphics[width=0.8\columnwidth]{usecase/UC3}
    \caption{UC3: Generazione assistita di unit test e E2E frontend}
  \end{figure}

  \usecaseactors{Sviluppatore, Agente AI}
  \usecasepre{L'ambiente di sviluppo \emph{Angular} del \emph{framework} è aperto, l'agente AI dispone delle \emph{skill} dedicate al \emph{testing frontend} e gli strumenti di \emph{unit test} e di test \emph{end-to-end} sono installati}
  \usecasedesc{Lo sviluppatore vuole introdurre \emph{unit test} sulla logica di un componente \emph{Angular} e/o test \emph{end-to-end} su un flusso utente, avvalendosi dell'agente AI}
  \usecasepost{Sono disponibili \emph{unit test} e test \emph{end-to-end} nella suite del \emph{frontend}}
  \usecasescenario{
    \begin{enumerate}
      \item Lo sviluppatore identifica il componente o il flusso utente (UC3.1);
      \item L'agente genera gli \emph{unit test} per la logica del componente (UC3.2);
      \item L'agente genera il test \emph{end-to-end} corrispondente (UC3.3);
      \item Lo sviluppatore revisiona e integra i test nella suite (UC3.4).
    \end{enumerate}
  }
  \label{uc:test-frontend}
\end{usecase}

\begin{usecase}{3.1}{Identificazione del componente o del flusso utente}
  \usecasedesc{Lo sviluppatore individua un componente da coprire con \emph{unit test} e/o un flusso utente da coprire con un test \emph{end-to-end}}
  \label{uc:test-frontend-target}
\end{usecase}

\begin{usecase}{3.2}{Generazione di unit test per la logica del componente}
  \usecasedesc{L'agente analizza la logica del componente \emph{Angular}, dei servizi iniettati e dei \emph{template binding}, e genera \emph{unit test} con i \emph{mock} delle dipendenze}
  \label{uc:test-frontend-unit}
\end{usecase}

\begin{usecase}{3.3}{Generazione del test end-to-end}
  \usecasedesc{L'agente genera un test \emph{end-to-end} che simula l'interazione dell'utente sull'applicazione, verificando le transizioni di stato e gli elementi attesi nelle varie schermate}
  \label{uc:test-frontend-e2e}
\end{usecase}

\begin{usecase}{3.4}{Revisione e integrazione nella suite}
  \usecasedesc{Lo sviluppatore esegue i test in locale, ne valuta l'esito e li integra nella suite \emph{frontend}}
  \label{uc:test-frontend-review}
\end{usecase}

\begin{usecase}{4}{Esecuzione della suite di test e produzione del report di copertura}
  \begin{figure}[!h]
    \centering
    \includegraphics[width=0.9\columnwidth]{usecase/UC4}
    \caption{UC4: Esecuzione della suite di test e produzione del report di copertura}
  \end{figure}

  \usecaseactors{Sviluppatore, Sia.Core}
  \usecasepre{La suite di test (\emph{unit}, \emph{integration} ed \emph{end-to-end}) è integrata nel processo di \emph{build}}
  \usecasedesc{Lo sviluppatore lancia l'esecuzione completa della suite per ottenere un esito complessivo e un report di copertura aggiornato del codice}
  \usecasepost{Lo sviluppatore dispone di un esito complessivo della suite e di un report di copertura}
  \usecasescenario{
    \begin{enumerate}
      \item Lo sviluppatore avvia la suite di test (UC4.1);
      \item Il sistema esegue tutti i test integrati nella \emph{build} (UC4.2);
      \item Il sistema produce il report di copertura del codice (UC4.3).
    \end{enumerate}
  }
  \usecaseextensions{
    \begin{enumerate}
      \item In presenza di test falliti, lo sviluppatore avvia la procedura di gestione (UC4.4).
    \end{enumerate}
  }
  \label{uc:esecuzione-test}
\end{usecase}

\begin{usecase}{4.1}{Avvio della suite di test}
  \usecasedesc{Lo sviluppatore lancia il comando o l'azione di avvio della suite}
  \label{uc:esecuzione-test-start}
\end{usecase}

\begin{usecase}{4.2}{Esecuzione dei test integrata nella build}
  \usecasedesc{Il sistema esegue in sequenza \emph{unit test}, \emph{integration test} ed \emph{end-to-end test} riportando l'esito di ciascuno}
  \label{uc:esecuzione-test-run}
\end{usecase}

\begin{usecase}{4.3}{Generazione del report di copertura del codice}
  \usecasedesc{Il sistema raccoglie i dati di copertura dei test e produce un report consultabile dallo sviluppatore}
  \label{uc:esecuzione-test-coverage}
\end{usecase}

\begin{usecase}{4.4}{Gestione dei test falliti}
  \usecasedesc{Lo sviluppatore consulta l'esito del test fallito, individua la causa nel codice sorgente e procede alla correzione del codice o del test}
  \label{uc:esecuzione-test-failure}
\end{usecase}

\begin{usecase}{5}{Propagazione della telemetria a un nuovo modulo}
  \begin{figure}[!h]
    \centering
    \includegraphics[width=0.9\columnwidth]{usecase/UC5}
    \caption{UC5: Propagazione della telemetria a un nuovo modulo}
  \end{figure}

  \usecaseactors{Sviluppatore, Agente AI}
  \usecasepre{Esiste un modulo del \emph{framework} non ancora coperto da telemetria, sono disponibili le \emph{skill} per la propagazione di \emph{trace}, metriche e \emph{log}, e lo stack \emph{Grafana LGTM} è raggiungibile}
  \usecasedesc{Lo sviluppatore vuole estendere la copertura della telemetria a un nuovo modulo del \emph{framework}, avvalendosi delle \emph{skill} dedicate dell'agente AI per garantire l'aderenza alle convenzioni}
  \usecasepost{Il modulo emette segnali di telemetria conformi alle convenzioni del \emph{framework} e visibili nello stack \emph{Grafana LGTM}}
  \usecasescenario{
    \begin{enumerate}
      \item Lo sviluppatore identifica il modulo da strumentare (UC5.1);
      \item Lo sviluppatore invoca la \emph{skill} di propagazione della telemetria (UC5.2);
      \item L'agente introduce il pattern \texttt{SiaXxxDiagnostics} e registra la sorgente nel \emph{pipeline} \gls{otel} (UC5.3);
      \item L'agente definisce metriche studiate per il modulo da coprire (UC5.4);
      \item Lo sviluppatore valida i segnali sullo stack \emph{Grafana LGTM} (UC5.5).
    \end{enumerate}
  }
  \label{uc:propagazione-telemetria}
\end{usecase}

\begin{usecase}{5.1}{Identificazione del modulo da strumentare}
  \usecasedesc{Lo sviluppatore individua un modulo non ancora coperto e ne valuta gli aspetti rilevanti per la telemetria (operazioni critiche, dipendenze esterne, prestazioni)}
  \label{uc:propagazione-target}
\end{usecase}

\begin{usecase}{5.2}{Invocazione della skill di propagazione della telemetria}
  \usecasedesc{Lo sviluppatore richiama l'agente AI passando il riferimento al componente, l'agente riconosce il contesto del \emph{prompt} e recupera la \emph{skill} di propagazione della telemetria}
  \label{uc:propagazione-invoke}
\end{usecase}

\begin{usecase}{5.3}{Introduzione della classe di diagnostica e registrazione in OTel}
  \usecasedesc{L'agente introduce nel modulo il pattern \texttt{SiaXxxDiagnostics}, definisce \texttt{ActivitySource} e \texttt{Meter} e registra la nuova sorgente nel \emph{pipeline} \gls{otel} del \emph{framework}}
  \label{uc:propagazione-diagnostics}
\end{usecase}

\newpage
\begin{usecase}{5.4}{Definizione di metriche studiate per il modulo da coprire}
  \usecasedesc{L'agente definisce le metriche del modulo applicando le convenzioni del \emph{framework} sui tag, in modo da evitare l'esplosione della cardinalità lato Prometheus}
  \label{uc:propagazione-metrics}
\end{usecase}

\begin{usecase}{5.5}{Validazione dei segnali sullo stack Grafana LGTM}
  \usecasedesc{Lo sviluppatore esegue il modulo in sviluppo e verifica che \emph{trace}, metriche e \emph{log} compaiano correttamente in \emph{Tempo}, \emph{Prometheus} e \emph{Loki}}
  \label{uc:propagazione-validate}
\end{usecase}

\begin{usecase}{6}{Visualizzazione di metriche, trace e log di un modulo}
  \begin{figure}[!h]
    \centering
    \includegraphics[width=0.9\columnwidth]{usecase/UC6}
    \caption{UC6: Visualizzazione di metriche, trace e log di un modulo}
  \end{figure}

  \usecaseactors{Sviluppatore}
  \usecasepre{Lo stack \emph{Grafana LGTM} è raggiungibile, il modulo di interesse è strumentato e produce segnali, e la dashboard del modulo è disponibile}
  \usecasedesc{Lo sviluppatore consulta in modo correlato metriche, \emph{trace} e \emph{log} di un modulo per ottenerne una visione unificata in un periodo di tempo}
  \usecasepost{Lo sviluppatore dispone di una visione unificata del comportamento del modulo nel periodo selezionato}
  \usecasescenario{
    \begin{enumerate}
      \item Lo sviluppatore apre la dashboard del modulo (UC6.1);
      \item Lo sviluppatore seleziona servizio, installazione e cliente tramite le variabili \emph{template} (UC6.2);
      \item Lo sviluppatore consulta in modo correlato metriche, \emph{trace} e \emph{log} (UC6.3).
    \end{enumerate}
  }
  \label{uc:dashboard-modulo}
\end{usecase}
\begin{usecase}{6.1}{Apertura della dashboard del modulo}
  \usecasedesc{Lo sviluppatore apre la dashboard associata al modulo dalla cartella di organizzazione del \emph{framework}}
  \label{uc:dashboard-modulo-open}
\end{usecase}

\begin{usecase}{6.2}{Selezione delle variabili template}
  \usecasedesc{Lo sviluppatore seleziona dai \emph{dropdown} delle variabili \emph{template} servizio, installazione e cliente per circoscrivere la visualizzazione}
  \label{uc:dashboard-modulo-filter}
\end{usecase}

\begin{usecase}{6.3}{Consultazione correlata di metriche, trace e log}
  \usecasedesc{Lo sviluppatore osserva i pannelli di metriche, \emph{trace} e \emph{log} relativi al modulo nel periodo selezionato, sfruttando i collegamenti reciproci messi a disposizione dalla dashboard}
  \label{uc:dashboard-modulo-correlate}
\end{usecase}

\begin{usecase}{7}{Diagnosi di una richiesta lenta tramite exemplar}
  \begin{figure}[!h]
    \centering
    \includegraphics[width=0.9\columnwidth]{usecase/UC7}
    \caption{UC7: Diagnosi di una richiesta lenta tramite exemplar}
  \end{figure}

  \usecaseactors{Sviluppatore}
  \usecasepre{Una dashboard \emph{Grafana} mostra un pannello percentili (es.\ p95/p99) di una metrica di tipo \emph{Histogram} con \emph{exemplar} abilitati e i \emph{log} sono indicizzati in \emph{Loki} con \texttt{TraceId}}
  \usecasedesc{Lo sviluppatore vuole ricostruire la catena diagnostica completa di una richiesta lenta a partire da un campione anomalo sul pannello percentili}
  \usecasepost{Lo sviluppatore ha ricostruito la catena diagnostica della singola richiesta lenta, dall'aggregato al \emph{log} di dettaglio}
  \usecasescenario{
    \begin{enumerate}
      \item Lo sviluppatore individua il campione anomalo sul pannello percentili (UC7.1);
      \item Lo sviluppatore naviga all'\emph{exemplar} associato (UC7.2);
      \item Lo sviluppatore apre la \emph{trace} corrispondente in \emph{Tempo} (UC7.3);
      \item Lo sviluppatore naviga ai \emph{log} correlati tramite \texttt{TraceId} (UC7.4).
    \end{enumerate}
  }
  \label{uc:diagnosi-exemplar}
\end{usecase}

\begin{usecase}{7.1}{Individuazione del campione anomalo}
  \usecasedesc{Lo sviluppatore identifica il punto temporale e il valore del campione anomalo sul pannello}
  \label{uc:diagnosi-exemplar-anomaly}
\end{usecase}

\begin{usecase}{7.2}{Navigazione all'exemplar associato}
  \usecasedesc{Lo sviluppatore clicca sull'\emph{exemplar} associato al campione, rappresentato come marker sul pannello}
  \label{uc:diagnosi-exemplar-link}
\end{usecase}

\begin{usecase}{7.3}{Apertura della trace in Tempo}
  \usecasedesc{Lo sviluppatore apre la \emph{trace} in \emph{Tempo}, esamina lo \emph{span} radice e i suoi figli, individuando il punto in cui si concentra la latenza}
  \label{uc:diagnosi-exemplar-trace}
\end{usecase}

\begin{usecase}{7.4}{Navigazione ai log correlati tramite TraceId}
  \usecasedesc{Lo sviluppatore segue il link verso \emph{Loki}, filtrato per \texttt{TraceId}, e consulta i \emph{log} emessi nel contesto della richiesta}
  \label{uc:diagnosi-exemplar-logs}
\end{usecase}

\newpage
\begin{usecase}{8}{Identificazione dell'utente coinvolto in una trace}
  \begin{figure}[!h]
    \centering
    \includegraphics[width=0.9\columnwidth]{usecase/UC8}
    \caption{UC8: Identificazione dell'utente coinvolto in una trace}
  \end{figure}

  \usecaseactors{Sviluppatore}
  \usecasepre{La telemetria del \emph{framework} arricchisce automaticamente ogni \emph{span} con gli attributi \texttt{enduser.id} e \texttt{enduser.name}, e \emph{log} e \emph{trace} sono indicizzati su tali attributi}
  \usecasedesc{Lo sviluppatore vuole risalire all'utente che ha originato una richiesta tracciata e analizzare il suo intero contesto}
  \usecasepost{Lo sviluppatore conosce l'identità dell'utente coinvolto e può proseguire l'indagine su tutto il suo contesto}
  \usecasescenario{
    \begin{enumerate}
      \item Lo sviluppatore apre la \emph{trace} di interesse (UC8.1);
      \item Lo sviluppatore legge gli attributi \texttt{enduser.id} e \texttt{enduser.name} dello \emph{span} radice (UC8.2);
      \item Lo sviluppatore filtra ulteriori \emph{trace} e \emph{log} per lo stesso utente (UC8.3).
    \end{enumerate}
  }
  \label{uc:enduser-trace}
\end{usecase}

\begin{usecase}{8.1}{Apertura della trace di interesse}
  \usecasedesc{Lo sviluppatore apre la \emph{trace} in \emph{Tempo}}
  \label{uc:enduser-trace-open}
\end{usecase}

\begin{usecase}{8.2}{Lettura degli attributi enduser dello span radice}
  \usecasedesc{Lo sviluppatore esamina gli attributi dello \emph{span} radice e legge \texttt{enduser.id} e \texttt{enduser.name}}
  \label{uc:enduser-trace-read}
\end{usecase}

\begin{usecase}{8.3}{Filtro di ulteriori trace e log per lo stesso utente}
  \usecasedesc{Lo sviluppatore esegue ricerche in \emph{Tempo} e \emph{Loki} filtrando per \texttt{enduser.id}, ottenendo l'intero contesto delle attività dell'utente}
  \label{uc:enduser-trace-filter}
\end{usecase}

\begin{usecase}{9}{Creazione o modifica di una dashboard tramite agente AI e MCP}
  \begin{figure}[!h]
    \centering
    \includegraphics[width=0.9\columnwidth]{usecase/UC9}
    \caption{UC9: Creazione o modifica di una dashboard tramite agente AI e MCP}
  \end{figure}

  \usecaseactors{Sviluppatore, Agente AI}
  \usecasepre{Lo stack \emph{Grafana LGTM} è raggiungibile e l'agente AI dispone delle credenziali \gls{mcp} verso \emph{Grafana} e verso le sorgenti dati}
  \usecasedesc{Lo sviluppatore vuole creare o modificare una dashboard delegando all'agente AI sia la formulazione delle query sia la pubblicazione finale}
  \usecasepost{La dashboard è disponibile in \emph{Grafana} con i pannelli e le query richiesti, validati sui dati reali}
  \usecasescenario{
    \begin{enumerate}
      \item Lo sviluppatore specifica all'agente le esigenze di \emph{monitoring} (UC9.1);
      \item L'agente formula e valida le query via \gls{mcp} (UC9.2);
      \item L'agente pubblica o aggiorna la dashboard sull'istanza \emph{Grafana} (UC9.3).
    \end{enumerate}
  }
  \usecaseextensions{
    \begin{enumerate}
      \item In presenza di una query non valida, l'agente avvia la procedura di gestione (UC9.4).
    \end{enumerate}
  }
  \label{uc:dashboard-mcp}
\end{usecase}

\begin{usecase}{9.1}{Specifica delle esigenze di monitoring}
  \usecasedesc{Lo sviluppatore descrive in linguaggio naturale i pannelli desiderati, le grandezze da osservare e i filtri di interesse}
  \label{uc:dashboard-mcp-spec}
\end{usecase}

\begin{usecase}{9.2}{Validazione delle query PromQL/TraceQL/LogQL via MCP}
  \usecasedesc{L'agente costruisce le query PromQL/TraceQL/LogQL necessarie e le esegue via \gls{mcp} contro le sorgenti dati per verificarne la sintassi e la pertinenza dei risultati}
  \label{uc:dashboard-mcp-query}
\end{usecase}

\begin{usecase}{9.3}{Pubblicazione o aggiornamento della dashboard}
  \usecasedesc{L'agente, tramite \gls{mcp}, crea o aggiorna la dashboard direttamente sull'istanza \emph{Grafana}, applicando convenzioni di naming e di organizzazione in cartelle}
  \label{uc:dashboard-mcp-publish}
\end{usecase}

\begin{usecase}{9.4}{Gestione di una query non valida}
  \usecasedesc{L'agente analizza l'errore o l'anomalia, riformula la query e ritenta; in caso di esito negativo dopo più tentativi, riporta il problema allo sviluppatore richiedendo chiarimenti sulla specifica}
  \label{uc:dashboard-mcp-error}
\end{usecase}

\section{Requisiti}
\label{sec:codice-requisiti}

\subsection{Requisiti funzionali}

\begin{enumerate}[label=RF\arabic*.]
  \item Il sistema deve tracciare il percorso completo di ogni richiesta HTTP in
        ingresso, incluse le operazioni di accesso al database.
  \item Le operazioni \gls{crud} verso il database devono essere riconoscibili tramite
        attributi che espongono il metodo e la \emph{repository} da cui è stata
        scatenata la richiesta.
  \item Il sistema deve esporre metriche operative per i moduli centrali del
        \emph{framework} (autenticazione, componenti \emph{Angular} come griglie e
        \emph{form}, prestazioni hardware, accesso al database).
  \item I \emph{log} devono essere correlabili alle \emph{trace} tramite
        \texttt{TraceId} e \texttt{SpanId}.
  \item Il sistema deve essere attivabile e disattivabile tramite configurazione,
        senza richiedere modifiche al codice.
  \item Deve essere possibile configurare indipendentemente il campionamento per
        ciascuna sorgente di \emph{trace}.
  \item Deve essere possibile visualizzare \emph{trace} e metriche raggruppate in
        dashboard \emph{Grafana} costruite sullo stack \emph{LGTM}.
  \item Le dashboard di \emph{Grafana} devono raccogliere i dati da più servizi e
        installazioni di clienti diversi, e permettere di filtrare per gli stessi.
  \item Il sistema deve arricchire automaticamente ogni \emph{span} con l'identità
        dell'utente autenticato (\texttt{enduser.id}, \texttt{enduser.name}) senza
        richiedere modifiche ai singoli moduli applicativi.
  \item Le metriche di tipo \emph{Histogram} devono esporre \emph{exemplar}, ovvero
        puntatori al \texttt{TraceId} del campione, per abilitare la navigazione dalla
        metrica aggregata alla \emph{trace} corrispondente.
  \item Devono essere prodotti \emph{unit test} per la logica di business dei moduli
        core del \emph{framework backend}, eseguibili in modo automatizzato.
  \item Devono essere prodotti \emph{integration test} per la verifica
        dell'attendibilità dei dati presenti nel database, delle procedure SQL e delle
        prospettive di configurazione dei componenti.
  \item Devono essere prodotti \emph{unit test} lato \emph{frontend} per la verifica
        della logica applicativa dei componenti \emph{Angular}.
  \item Devono essere prodotti test \emph{end-to-end} lato \emph{frontend} per
        verificare il corretto comportamento e la corretta visualizzazione
        dell'applicazione dal punto di vista dell'utente finale.
  \item La suite di test deve essere integrabile nel processo di \emph{build} e
        fornire un report di copertura del codice.
  \item Devono essere realizzati agenti AI specializzati e \emph{skill} procedurali
        per propagare in modo uniforme telemetria e test ai diversi moduli del
        \emph{framework}.
\end{enumerate}

\subsection{Requisiti non funzionali}

\begin{enumerate}[label=RNF\arabic*.]
  \item La telemetria non deve degradare sensibilmente le prestazioni in produzione.
  \item Le metriche devono essere progettate per non generare un numero eccessivo di serie temporali in Prometheus; in particolare, i tag ad alta variabilità (es.\ identificativi utente o di record) non devono essere usati come attributi di metriche.
  \item I segnali prodotti devono essere compatibili con qualsiasi backend \gls{otel}-compatibile.
  \item Le impostazioni devono essere modificabili per ambiente (sviluppo, \emph{staging}, produzione) tramite \emph{appsettings}.
\end{enumerate}

\subsection{Tracciamento dei requisiti}

Dall'analisi dei requisiti e degli \emph{use case} è stata sviluppata una
tabella di tracciamento che pone in relazione ciascun requisito con i casi
d'uso che ne dipendono. Per distinguere le diverse classi di requisito si è
adottato un codice identificativo nella forma R(F/Q/V)(O/D), dove:

\begin{itemize}
  \item la prima lettera (sempre R) indica che si tratta di un requisito;
  \item la seconda lettera ne indica la natura: F (funzionale), Q
        (qualitativo) o V (di vincolo);
  \item la terza lettera ne indica la priorità: O (obbligatorio) o D
        (desiderabile).
\end{itemize}

Per non appesantire il corpo del capitolo, le tabelle
complete del tracciamento sono raccolte nell'appendice~\ref{app:requisiti},
alle quali si rimanda i lettori interessati al dettaglio.
